\section{Conclusion}
\label{sec:conclusion}

This study presented a lightweight and interpretable machine learning framework for gait symmetry assessment using bilateral gyroscope data from shank-mounted IMUs. 
By segmenting raw sensor signals into individual gait cycles and extracting a minimal set of meaningful features—including stance time, swing time, 
stride duration, and motion dynamics—we demonstrated that reliable classification of gait asymmetry can be achieved without relying on complex or computationally 
intensive methods.

The introduction of a three-tier confidence labeling scheme enabled more nuanced training and evaluation, accounting for varying degrees of asymmetry severity. 
Classifiers such as KNN, XGBoost, and Random Forest consistently outperformed other models, achieving high accuracy and balanced precision-recall performance 
across all confidence levels. Furthermore, recursive feature selection confirmed that a small subset of time-domain features is sufficient for robust model 
performance, reinforcing the feasibility of real-time deployment on mobile or embedded platforms.

These findings highlight the potential of combining wearable sensor data with transparent machine learning techniques to support stroke rehabilitation, remote 
monitoring, and adaptive feedback systems. Future work will involve validating the system in real-time settings and expanding the dataset to include broader 
pathological gait conditions for further generalization.